%%%=============================================================================
%%%  MANUAL DEL ESTUDIANTE --- Tracker UAI para el Laboratorio 1
%%%  Instalación y uso paso a paso (primera vez)
%%%  Universidad Adolfo Ibáñez / FACH --- Prof. David Aguayo Vera
%%%
%%%  Requiere  uaiestilo.sty  en la misma carpeta.
%%%  Compilar con pdfLaTeX (dos pasadas para índice y referencias).
%%%=============================================================================
\documentclass[11pt,a4paper]{article}

\usepackage{uaiestilo}

%%-----------------------------------------------------------------------------
%% Marcadores de acción (para distinguir qué hace el estudiante)
%%-----------------------------------------------------------------------------
\newcommand{\Accion}{\textcolor{uaiAzul}{\faHandPointRight\ \textbf{Acción:}}\ }
\newcommand{\Ingresa}{\textcolor{uaiNaranja}{\faKeyboard\ \textbf{Ingresa:}}\ }
\newcommand{\Verifica}{\textcolor{uaiVerde}{\faCheckCircle\ \textbf{Verifica:}}\ }
\newcommand{\SiFalla}{\textcolor{uaiRojo}{\faExclamationTriangle\ \textbf{Si falla:}}\ }
\newcommand{\Guarda}{\textcolor{uaiVioleta}{\faSave\ \textbf{Guarda:}}\ }

%% Caja para indicar dónde conviene poner una captura de pantalla
\newtcolorbox{captura}[1][]{%
  enhanced, breakable, colback=uaiNiebla, colframe=uaiAcero,
  boxrule=0.6pt, arc=1.2mm, left=8pt, right=8pt, top=6pt, bottom=6pt,
  fonttitle=\bfseries\color{uaiTinta},
  title={\faCamera\ Captura sugerida\ifx\relax#1\relax\else\ ---\ #1\fi}}

%%-----------------------------------------------------------------------------
%% Metadatos
%%-----------------------------------------------------------------------------
\curso{Universidad Adolfo Ibáñez --- Facultad de Ingeniería y Ciencias}
\cursocorto{Tracker UAI --- Manual del estudiante (Lab 1)}
\title{Manual del estudiante: Tracker UAI}
\subtitle{Instalación y uso paso a paso para el Laboratorio 1}
\author{Prof.\ David Aguayo Vera}
\date{Segundo semestre de 2026}

\uailogos{%
  \uailogo[1.1cm]{logo-uai.png}\hfill%
  \uailogo[1.1cm]{logo-fach.png}}

%%=============================================================================
\begin{document}
%%=============================================================================

\maketitle

\begin{obs}[Cómo leer este manual]
  Está pensado para usar el programa \textbf{por primera vez}. Sigue los pasos
  en orden. A lo largo del texto verás marcadores de color que indican qué
  hacer en cada momento:
  \begin{itemize}
    \item \Accion algo que debes hacer (un clic, arrastrar, navegar).
    \item \Ingresa un dato que debes escribir.
    \item \Verifica un resultado que debes comprobar en pantalla.
    \item \SiFalla qué hacer si no sale lo esperado.
    \item \Guarda un archivo que debes conservar para tu informe.
  \end{itemize}
  Las cajas grises con el ícono de cámara indican dónde conviene mirar (o pegar)
  una captura de pantalla.
\end{obs}

\tableofcontents
\vspace{4mm}

%%=============================================================================
\section{Introducción}
%%=============================================================================

\textbf{Tracker UAI} es un programa para \textbf{analizar videos de experimentos
de física}. Se abre en tu navegador web y te guía por etapas para medir cómo se
mueve un objeto en un video y obtener resultados numéricos.

\subsection{¿Para qué se usa en el Laboratorio 1?}

En el Laboratorio~1 grabarás un experimento con tu celular (por ejemplo, un carro
que desliza por un plano inclinado) y usarás el programa para \textbf{seguir la
posición del objeto cuadro a cuadro}. Con esos datos el programa calcula la
velocidad y la aceleración, y ajusta un modelo físico para entregarte, por
ejemplo, la \textbf{aceleración del carro} en cada lanzamiento.

\subsection{¿Qué información permite obtener?}

\begin{idea}
  A partir del video, el programa entrega:
  \begin{itemize}
    \item La \textbf{posición} del objeto en el tiempo, en metros.
    \item Su \textbf{velocidad} y \textbf{aceleración}.
    \item Gráficos de posición, velocidad, aceleración y trayectoria.
    \item El \textbf{ajuste de un modelo} (por ejemplo, la aceleración $a$) con su
      \textbf{incertidumbre} y el coeficiente $R^2$.
    \item Una tabla de datos y figuras que puedes \textbf{exportar} para tu informe.
  \end{itemize}
\end{idea}

Además, el programa incluye una herramienta de \emph{Gráficos} para dibujar y
ajustar datos escritos a mano (útil para el gráfico final del grupo, por ejemplo
$a$ en función de $\sin\theta$), y una de \emph{Medición} para medir distancias y
ángulos sobre una foto.

%%=============================================================================
\section{Requisitos previos}
%%=============================================================================

\subsection{Requisitos del computador}
\begin{itemize}
  \item Windows, macOS o Linux (cualquiera de los tres sirve).
  \item Conexión a internet \textbf{solo la primera vez} (para instalar).
  \item Un navegador web moderno (Chrome, Edge, Firefox o Safari).
\end{itemize}

\subsection{Programa que debes tener instalado}
\begin{itemize}
  \item \textbf{Python 3.11 o superior.} Es lo único que instalas tú; el resto lo
    instala el programa solo.
\end{itemize}

\subsection{Archivos que debes descargar antes de comenzar}
\begin{itemize}
  \item La \textbf{carpeta del programa} (Tracker UAI), que te entregará el
    profesor o descargarás desde el repositorio del curso.
  \item Tu(s) \textbf{video(s)} del experimento, grabados con el celular.
\end{itemize}

%%=============================================================================
\section{Instalación del programa}
%%=============================================================================

\subsection{Paso 1: instalar Python (una sola vez)}

\begin{receta}[Instalar Python según tu sistema]
  \begin{itemize}
    \item \textbf{Windows:} entra a \texttt{python.org/downloads}, descarga
      Python 3.11 o superior y, en el instalador, \textbf{marca la casilla
      «Add Python to PATH»} antes de continuar.
    \item \textbf{macOS:} descarga desde \texttt{python.org/downloads}, o con
      Homebrew ejecuta \texttt{brew install python}.
    \item \textbf{Linux:} instala \texttt{python3}, \texttt{python3-venv} y
      \texttt{python3-pip} con el gestor de tu distribución.
  \end{itemize}
\end{receta}

\begin{errcomun}
  En Windows, \textbf{olvidar marcar «Add Python to PATH»} es el error más común:
  después el programa no encuentra Python y no abre. Si te pasa, vuelve a instalar
  Python marcando esa casilla.
\end{errcomun}

\subsection{Paso 2: ubicar la carpeta del programa}

\Accion Descomprime la carpeta del programa (si vino en un \texttt{.zip}: clic
derecho \textrightarrow{} \emph{Extraer todo}) en un lugar fácil de encontrar,
por ejemplo el Escritorio o Descargas. Dentro verás archivos como
\texttt{app.py}, \texttt{requirements.txt} y los lanzadores.

\subsection{Paso 3: ejecutar el programa por primera vez}

\begin{receta}[Abrir Tracker UAI]
  \begin{itemize}
    \item \textbf{Windows:} haz \textbf{doble clic} en
      \texttt{Iniciar\_Tracker\_UAI.bat}.
    \item \textbf{macOS / Linux:} abre una terminal en la carpeta y ejecuta
      \texttt{bash iniciar\_tracker.sh}.
  \end{itemize}
\end{receta}

\begin{captura}[Ventana negra (terminal) mostrando la instalación por primera vez]
  La primera vez, la ventana de la terminal mostrará el avance de la instalación.
\end{captura}

\Verifica La \textbf{primera vez} tarda algunos minutos porque instala lo
necesario; la ventana negra queda abierta mostrando el avance. Luego el programa
se abre solo en tu navegador (dirección \texttt{http://localhost:8501}). Las
veces siguientes abre en pocos segundos.

\subsection{Si Windows muestra una advertencia de seguridad}

\begin{alerta}
  Si aparece una pantalla azul que dice \emph{«Windows protegió tu PC»}, es una
  advertencia normal para archivos descargados. Haz clic en \textbf{«Más
  información»} y luego en \textbf{«Ejecutar de todas formas»}.
\end{alerta}

\subsection{Método alternativo (por terminal)}

Si prefieres no usar el lanzador, en una terminal dentro de la carpeta:
\begin{lstlisting}[caption={Instalar y ejecutar manualmente.}]
pip install -r requirements.txt
python -m streamlit run app.py
\end{lstlisting}

\begin{errcomun}
  Si al abrir sale un error que menciona \texttt{ImportError} o
  \texttt{UseColumnWith}, tienes versiones incompatibles de las librerías.
  Solución: ejecuta \texttt{pip install -r requirements.txt} de nuevo (el archivo
  ya fija las versiones correctas).
\end{errcomun}

%%=============================================================================
\section{Descripción general de la interfaz}
%%=============================================================================

Al abrir el programa verás dos zonas: la \textbf{barra lateral} (a la izquierda) y
el \textbf{área principal} (al centro).

\begin{captura}[Vista general con la barra lateral y el área principal]
\end{captura}

\subsection{Barra lateral (izquierda)}
\begin{itemize}
  \item \textbf{Sube tu video}: botón para cargar el archivo de video.
  \item \textbf{Etapa}: lista para moverte entre las etapas del análisis
    (1~Video, 2~Calibración, 3~Tracking, 4~Cinemática, 5~Ajuste, 6~Exportar),
    más las herramientas \emph{Gráficos (datos manuales)} y \emph{Medición (foto)}.
  \item \textbf{Opciones avanzadas}: aquí se activa el fps de captura manual
    (para videos en cámara lenta).
\end{itemize}

\subsection{Área principal (centro)}
Muestra la etapa seleccionada. Según la etapa, verás el video, controles de
navegación, la imagen para hacer clic, tablas de datos y gráficos.

\begin{idea}
  El trabajo avanza \textbf{de arriba hacia abajo} y \textbf{de la Etapa 1 a la
  6}. Cambiar de etapa \emph{no} borra tu trabajo: la calibración y los puntos
  marcados se conservan.
\end{idea}

\subsection{Controles que verás a menudo}
\begin{itemize}
  \item Botones de navegación de frames: \texttt{-10}, \texttt{-5}, \texttt{-1},
    \texttt{+1}, \texttt{+5}, \texttt{+10}.
  \item Campo \textbf{«Ir al frame»}: para escribir el número de frame exacto.
  \item Barra deslizante (\emph{slider}) para moverte por el video.
  \item La \textbf{imagen del video}: al hacer clic sobre ella marcas puntos.
\end{itemize}

%%=============================================================================
\section{Preparación de los datos del laboratorio}
%%=============================================================================

\subsection{Qué archivos necesitas}
\begin{itemize}
  \item Uno o varios \textbf{videos} del experimento (uno por lanzamiento).
  \item Saber la \textbf{distancia real} de un objeto de referencia visible en el
    video (por ejemplo, el largo de una regla), para calibrar.
\end{itemize}

\subsection{Cómo organizar los archivos}
\Accion Crea una carpeta por experimento y nombra los videos de forma clara, por
ejemplo \texttt{lanzamiento\_10grados.mp4}, \texttt{lanzamiento\_15grados.mp4}, etc.

\subsection{Formato recomendado de los videos}
\begin{itemize}
  \item Formatos aceptados: \texttt{mp4}, \texttt{mov}, \texttt{avi},
    \texttt{mkv}, \texttt{m4v}, \texttt{webm}.
  \item Recomendado: \textbf{mp4 con códec H.264}.
  \item Duración corta (solo el tramo del movimiento).
\end{itemize}

\begin{estrategia}[Recomendaciones para evitar problemas]
  \begin{itemize}
    \item Graba con la cámara \textbf{perpendicular} al plano del movimiento y
      \textbf{fija} (trípode o apoyada).
    \item Coloca la \textbf{regla de calibración en el mismo plano} que el
      movimiento.
    \item Buena iluminación; el objeto bien visible y, si es posible, de un
      \textbf{color distinto al fondo} (facilita la detección automática).
    \item Recorta el video para que dure poco: menos frames, más rápido.
  \end{itemize}
\end{estrategia}

%%=============================================================================
\section{Realización del análisis paso a paso}
%%=============================================================================

Esta es la columna vertebral del trabajo. Sigue las etapas en orden.

\subsection{Etapa 1 --- Cargar el video y elegir el tramo}

\Accion En la barra lateral, pulsa \textbf{«Sube tu video»} y elige tu archivo.

\Verifica Arriba aparecen los datos del video (número de frames, fps, resolución,
duración). La primera vez, espera el mensaje \emph{«Preparando el video…»}.

\begin{captura}[Etapa 1 con el visor de frames y los datos del video]
\end{captura}

\Accion Con los botones (\texttt{-10}/\texttt{-5}/\texttt{-1}/\texttt{+1}/%
\texttt{+5}/\texttt{+10}), el campo \textbf{«Ir al frame»} y la barra, ubica el
\textbf{frame donde empieza} el movimiento y el \textbf{frame donde termina}.

\Accion Pulsa \textbf{«Usar frame actual como INICIO»} en el primero, y navega al
final y pulsa \textbf{«Usar frame actual como FINAL»}.

\Verifica Abajo aparece un mensaje verde: \emph{«Intervalo: frame X $\to$ Y…»}.

\begin{ojo}
  Si cambias el intervalo después de trackear, el programa \textbf{borra los
  puntos fuera del nuevo rango} y recalcula todo. Elige bien el intervalo antes
  de marcar puntos.
\end{ojo}

\subsection{Etapa 2 --- Calibración (ver Sección~\ref{sec:calib})}
\subsection{Etapa 3 --- Detección y seguimiento (ver Sección~\ref{sec:track})}
\subsection{Etapa 4 --- Cinemática (ver Sección~\ref{sec:proc})}
\subsection{Etapa 5 --- Ajuste del modelo (ver Sección~\ref{sec:proc})}
\subsection{Etapa 6 --- Exportar (ver Sección~\ref{sec:export})}

\begin{obs}
  Las etapas 2 a 6 se explican en detalle en las secciones siguientes, porque
  cada una tiene sus propios controles.
\end{obs}

%%=============================================================================
\section{Detección y seguimiento de los objetos}
\label{sec:track}
%%=============================================================================

En la \textbf{Etapa 3 (Tracking)} marcas la posición del objeto en cada frame.
Hay dos modos: \textbf{Manual} y \textbf{Automático}.

\subsection{Navegación por frames}
\begin{itemize}
  \item \Accion Botones \texttt{-1}/\texttt{+1} para avanzar un frame; \texttt{-5}%
    /\texttt{+5} y \texttt{-10}/\texttt{+10} para saltos rápidos en videos largos.
  \item \Accion Campo \textbf{«Ir al frame»}: escribe el número exacto y salta ahí.
\end{itemize}

\subsection{Seleccionar el intervalo a analizar}
Se hace en la Etapa~1 (ver arriba). En la Etapa~3, el rango de detección
automática viene \textbf{pre-cargado con ese intervalo}.

\subsection{Detección automática}
\Accion Elige el modo \textbf{Automático} y el método:
\begin{itemize}
  \item \textbf{Plantilla}: \Accion arrastra un recuadro sobre el objeto en el
    frame inicial.
  \item \textbf{Color}: \Accion haz clic sobre el objeto para tomar su color;
    ajusta la tolerancia y observa la \textbf{máscara} (blanco = detectado).
  \item \textbf{CSRT}: aparece solo si está disponible; también se arrastra un
    recuadro.
\end{itemize}

\Accion Pulsa \textbf{«Ejecutar seguimiento»}.

\Verifica Verás una \textbf{vista previa en vivo} y, al terminar, la imagen con
\textbf{todas las detecciones superpuestas}.

\begin{captura}[Detección automática con los puntos superpuestos]
\end{captura}

\SiFalla Si no detecta bien (puntos saltando, objeto perdido): cambia de método
(prueba \emph{Color} si el objeto tiene color distintivo), ajusta la tolerancia,
o usa el modo \textbf{Manual}.

\subsection{Revisar y corregir puntos}
\Accion En la sección \textbf{«Corrección frame por frame»} (o en el modo Manual),
navega frame a frame. Si un punto está mal, \Accion \textbf{haz clic en la
posición correcta} para moverlo; o pulsa \textbf{«Borrar punto de este frame»}.

\Verifica En la \textbf{tabla de datos} de abajo puedes revisar y editar los
valores; también eliminar filas.

\subsection{Marcado manual (alternativa)}
\Accion Elige modo \textbf{Manual}, fija «Avanzar cada N frames», y \Accion haz
\textbf{clic sobre el objeto} en cada frame: el video avanza solo.

%%=============================================================================
\section{Calibración y sistema de referencia}
\label{sec:calib}
%%=============================================================================

La \textbf{Etapa 2 (Calibración)} le enseña al programa a convertir píxeles en
metros y dónde está el origen de coordenadas. \textbf{Sin calibrar, los resultados
no salen en metros.}

\subsection{Escala espacial (píxel $\to$ metro)}
\Accion En el selector «¿Qué vas a marcar…?», con \textbf{Punto 1} activo, haz
clic en un extremo de tu regla; luego, con \textbf{Punto 2}, en el otro extremo.

\Ingresa La \textbf{distancia real} entre esos dos puntos, en metros (por ejemplo,
\num{0.30} m para una regla de 30~cm).

\begin{captura}[Calibración: dos puntos sobre la regla y el origen marcados]
\end{captura}

\subsection{Origen de coordenadas}
\Accion Selecciona \textbf{Origen} y haz clic donde quieras el punto $(0,0)$. El
eje $Y$ físico apunta hacia \textbf{arriba} (como en física).

\subsection{Ejes con orientación libre (opcional)}
\Accion Si tu eje $x$ no es horizontal (por ejemplo, quieres alinearlo con un
plano inclinado), activa \textbf{«Ejes con orientación libre»} y marca un punto
en la dirección $+x$.

\Accion Pulsa \textbf{«Aplicar calibración»}.

\Verifica Aparece un mensaje verde con la escala (\emph{m/px}) y el origen.

\begin{idea}
  \textbf{Qué significa cada cosa:} la \emph{escala} dice cuántos metros mide un
  píxel; el \emph{origen} es el punto desde donde se miden las posiciones; los
  \emph{ejes} definen las direcciones $x$ e $y$. Si calibras mal (distancia
  equivocada o regla en otro plano), \textbf{todos} los resultados en metros
  saldrán mal.
\end{idea}

\SiFalla Si la escala parece rara (velocidades enormes o minúsculas), revisa que
la distancia ingresada esté en \textbf{metros} y que hayas marcado bien los dos
extremos de la regla.

%%=============================================================================
\section{Procesamiento y análisis de los datos}
\label{sec:proc}
%%=============================================================================

\subsection{Qué hace el programa con tus datos}
Con los puntos marcados y la calibración, el programa calcula el \textbf{tiempo},
la \textbf{posición en metros}, y por diferencias finitas la \textbf{velocidad} y
la \textbf{aceleración}. Todo esto está en la \textbf{Etapa 4 (Cinemática)}.

\subsection{Revisar tablas y gráficos}
\Verifica En la Etapa~4 verás los gráficos de \textbf{posición}, \textbf{velocidad},
\textbf{aceleración} y \textbf{trayectoria}. Abre \textbf{«Tabla con derivados»}
para ver los números.

\begin{captura}[Etapa 4: gráficos de posición, velocidad, aceleración y trayectoria]
\end{captura}

\Accion Si los datos se ven ruidosos (sobre todo la aceleración), activa
\textbf{«Suavizar posición»} y ajusta la ventana.

\subsection{Ajuste del modelo y ecuación (Etapa 5)}
\Accion Ve a la \textbf{Etapa 5 (Ajuste)} y elige el \textbf{preset de tu
laboratorio}. Por ejemplo, para el plano inclinado: \emph{«FIS101 · Lab 1 ·
Gravedad en plano inclinado»}.

\Accion Elige la variable (posición $\to$ parábola, o velocidad $\to$ recta).

\Verifica El programa muestra la \textbf{aceleración $a$ con su incertidumbre}, la
curva ajustada sobre los datos y el \textbf{$R^2$} (mientras más cerca de 1, mejor
el ajuste).

\begin{captura}[Etapa 5: datos con la curva ajustada, el valor de $a$ y el $R^2$]
\end{captura}

\subsection{Gráfico final del grupo (herramienta «Gráficos»)}
Para el resultado del grupo (por ejemplo, $a$ en función de $\sin\theta$ de tus 6
ángulos):
\begin{itemize}
  \item \Accion Ve a \textbf{«Gráficos (datos manuales)»}.
  \item \Accion Elige la plantilla \emph{«FIS101 · g desde a vs sen $\theta$»}.
  \item \Ingresa en la tabla los pares $(\sin\theta,\ a)$ de cada ángulo.
  \item \Verifica La app hace el ajuste lineal y muestra \textbf{$g$} (la
    pendiente) con su incertidumbre y la comparación con el valor teórico.
\end{itemize}

\subsection{Cómo interpretar los gráficos}
\begin{idea}
  \begin{itemize}
    \item \textbf{Posición vs.\ tiempo}: en el plano inclinado debe curvarse
      (parábola).
    \item \textbf{Velocidad vs.\ tiempo}: debe ser aproximadamente una recta; su
      pendiente es la aceleración.
    \item \textbf{$R^2$ cercano a 1}: el modelo describe bien tus datos.
    \item Puntos muy dispersos o fuera de la curva suelen indicar un
      \textbf{punto mal marcado}: vuelve a la Etapa~3 y corrígelo.
  \end{itemize}
\end{idea}

%%=============================================================================
\section{Resultados y exportación}
\label{sec:export}
%%=============================================================================

En la \textbf{Etapa 6 (Exportar)}:

\Accion Marca con \textbf{casillas} qué figuras quieres (posición, velocidad,
aceleración, trayectoria, error de detección) y si quieres el \textbf{CSV} de datos.

\Accion Pulsa \textbf{«Generar archivos»} y luego \textbf{«Descargar ZIP»} (o
«Descargar solo el CSV»).

\begin{captura}[Etapa 6: casillas de figuras y botón de descarga]
\end{captura}

\Guarda Para tu informe, conserva:
\begin{itemize}
  \item El \textbf{ZIP} con las figuras (PNG) y el \textbf{CSV} de datos.
  \item Del gráfico del grupo, el \textbf{PNG/SVG} exportado con el valor de $g$.
  \item Una \textbf{captura} de la Etapa~5 con el valor de $a$, su incertidumbre y
    el $R^2$ de cada lanzamiento.
\end{itemize}

%%=============================================================================
\section{Errores frecuentes y solución de problemas}
%%=============================================================================

\begin{table}[htbp]
  \centering
  \caption{Problemas frecuentes y su solución.}
  \label{tab:problemas}
  \begin{tabular}{@{}p{5.4cm} p{8.1cm}@{}}
    \toprule
    \textbf{Problema} & \textbf{Solución} \\
    \midrule
    El programa no abre & Revisa que Python esté instalado con «Add to PATH»
      (Windows). Si ves \texttt{ImportError/UseColumnWith}, corre otra vez
      \texttt{pip install -r requirements.txt}. \\
    El video no carga & Verifica el formato (mp4/mov/avi/mkv). Si es \texttt{.mov}
      de iPhone, conviértelo a H.264 (ver abajo). \\
    El video se ve mal o «se congela» & Suele ser el códec HEVC/H.265.
      Conviértelo a H.264. Si es muy largo, recórtalo. \\
    No se detectan bien los objetos & Cambia de método (prueba \emph{Color});
      ajusta la tolerancia; o usa el modo \emph{Manual}. \\
    Los puntos detectados son incorrectos & Corrígelos: navega al frame y haz
      clic en la posición correcta, o edita la tabla. \\
    La calibración da resultados extraños & Revisa que la distancia esté en
      \textbf{metros} y que marcaste bien los extremos de la regla. \\
    El programa parece detenido & Espera: la primera precarga del video o la
      detección pueden tardar. Mira la barra de progreso. \\
    Los gráficos no aparecen & Necesitas \textbf{calibración aplicada} y al menos
      2 puntos \emph{dentro del intervalo}. \\
    No se pueden exportar los resultados & Primero marca puntos y calibra; luego
      pulsa «Generar archivos» antes de descargar. \\
    \bottomrule
  \end{tabular}
\end{table}

\begin{alerta}
  \textbf{Convertir un video a H.264} (si el \texttt{.mov} no carga), con la
  herramienta \texttt{ffmpeg} desde la terminal:
\end{alerta}
\begin{lstlisting}[caption={Convertir un video a H.264.}]
ffmpeg -i entrada.mov -c:v libx264 -pix_fmt yuv420p salida.mp4
\end{lstlisting}

%%=============================================================================
\section{Procedimiento recomendado para el Laboratorio 1}
%%=============================================================================

\begin{receta}[Flujo completo, de principio a fin]
  \begin{enumerate}
    \item \textbf{Abre} el programa (doble clic en el lanzador) y espera a que
      cargue en el navegador.
    \item \textbf{Etapa 1:} sube el video, revísalo frame a frame y define el
      \textbf{intervalo} (frame inicial y final del movimiento).
    \item \textbf{Etapa 2:} calibra --- marca los dos puntos de la regla, ingresa
      la distancia real en metros, marca el origen y \textbf{aplica}.
    \item \textbf{Etapa 3:} trackea el objeto (automático o manual) y
      \textbf{revisa/corrige} los puntos.
    \item \textbf{Etapa 4:} mira los gráficos de posición, velocidad y
      aceleración; activa suavizado si hay ruido.
    \item \textbf{Etapa 5:} elige el preset de tu laboratorio y obtén el valor
      (por ejemplo, la aceleración $a$) con su incertidumbre y $R^2$. Anótalo.
    \item \textbf{Repite} las etapas 1--5 para cada lanzamiento/ángulo.
    \item \textbf{Gráficos:} con la plantilla de tu laboratorio, ingresa los pares
      del grupo (por ejemplo $\sin\theta$ vs $a$) y obtén el resultado final
      (por ejemplo $g$).
    \item \textbf{Etapa 6:} exporta las figuras y el CSV. \Guarda todos los
      archivos para el informe.
  \end{enumerate}
\end{receta}

\begin{resumen}
  \begin{itemize}
    \item El análisis va \textbf{en orden}, de la Etapa~1 a la~6.
    \item \textbf{Calibrar bien} es lo más importante: define los resultados en
      metros.
    \item Un \textbf{punto mal marcado} se nota como un dato fuera de la curva;
      corrígelo en la Etapa~3.
    \item Guarda \textbf{CSV + figuras + capturas} de los valores ajustados para
      el informe.
  \end{itemize}
\end{resumen}

\end{document}
